\section{Codiseño extremo: de la optimización de componentes al óptimo del sistema completo}

\subsection{Por qué es necesaria la optimización entre capas}

Jensen divide «extreme co-design» en dos capas. La primera es la coordinación entre software y hardware, desde architecture, chips y systems hasta system software, algorithms y applications. La segunda es la coordinación del sistema físico, que integra power y cooling en una misma tabla de presupuesto de desempeño. Un centro de datos de IA no es un problema puramente de software, ni puramente de chips, sino un problema de acoplamiento interdisciplinario.

\begin{importantbox}{El óptimo del sistema no es igual a la suma de los óptimos locales}
En el entrenamiento a gran escala, que un solo componente sea el más fuerte de manera aislada puede provocar que el sistema en su conjunto se degrade. Por ejemplo, una densidad de cómputo más agresiva, si desencadena fluctuaciones en el suministro eléctrico, limitaciones térmicas y congestión de red, terminará por reducir el token throughput. La formulación de Jensen es: tratar cada capa como una restricción y no como una premisa.
\end{importantbox}

\subsection{El límite de ingeniería que representa NVLink 72}

En la entrevista se menciona varias veces NVLink 72 y rack-scale computing. Lo importante no es solo la cantidad de GPU, sino que «la máquina completa pueda operar»: la densidad del cableado, la topología de conmutación, la mecánica del rack, los canales térmicos, las rutas de tolerancia a fallas y el proceso de pruebas de fábrica deben reescribirse por completo. Jensen enfatiza en particular que, a esta escala, la supercomputadora debe ensamblarse y validarse dentro de la cadena de suministro, y no transportarse al centro de datos para armarse ahí.

\begin{knowledgebox}{El hilo evolutivo de DGX a rack}
En la conversación aparecen palabras clave como DGX1, Grace Blackwell, Vera Rubin y NVLink72. Todas apuntan en conjunto a una tendencia: la forma del producto pasa de la «tarjeta» a la «máquina completa», y de ahí a la «unidad replicable de centro de datos». Esto también es la evidencia técnica de la transformación de NVIDIA de fabricante de componentes a «AI infrastructure company».
\end{knowledgebox}

\begin{warningbox}{Malentendido común: entender la infraestructura de IA como compra de GPU}
Si solo se observa la cantidad de GPU, se subestima el impacto de la red, el suministro eléctrico, la disipación de calor, el empaquetado y las pruebas sobre el ritmo de entrega. En los proyectos reales, lo que primero retrasa la puesta en marcha suele ser la remodelación de las salas de cómputo, la aprobación del suministro eléctrico, la validación del clúster y la automatización de operaciones, y no el envío de los chips en sí.
\end{warningbox}

\subsection{Principios de diseño en sistemas complejos}

Hacia la parte final, Jensen ofrece una frase clave: \textit{``be as complex as necessary, but as simple as possible.''} Esta frase puede entenderse como dos restricciones:
\begin{itemize}
    \item La complejidad no debe idealizarse: cualquier grado de complejidad debe tener un retorno en desempeño o confiabilidad.
    \item La simplificación no debe sacrificar el objetivo: hay que aceptar la complejidad donde sea necesaria y forzar la abstracción donde sea posible.
\end{itemize}

\subsection{Resumen de la sección}

La esencia de «extreme co-design» no es reunir a más expertos en una sola junta, sino tratar el cómputo, la red, la energía y la fabricación como un mismo problema de optimización. Una forma de producto como NVLink72 muestra que NVIDIA ya ha llevado la dimensión de la competencia de los «parámetros del chip» a la «capacidad de entrega del sistema».

